Remove row 73-82 and last sentence in row 83.\\

The parameters used to train the ML models are temperature, humidity, pressure and K96 sensor output. The training data used include values of the parameters varying between the maximum and minimum value that each of them shall stay in according to the experiment design and its requirements to ensure that the models are exposed to all relevant environmental conditions. \\
The calibration procedure to collect training and testing data will be conducted at the Norunda weather station, which provides the necessary infrastructure for controlled calibration. Norunda offers four different calibration gases with specified concentrations of $CH_4$ and $CO_2$, nitrogen gas, and a dew point generator for humidity variation. Norunda also posses a Los Gatos Research sensor that can be used for high accuracy reference measurements of atmospheric $CH_4$ and $CO_2$ concentrations, and thereby collect atmospheric test data to validate the calibration model later on.
% Unfortunately, it can't be used for the calibration itself because it is used for continuous atmospheric monitoring. However, the measurement setup can be connected to it via a tube and thereby collect atmospheric test data to validate the calibration model later on. 
Therefore, the simulation of the environmental conditions expected during the stratospheric flight and the collection of reliable reference data for testing is possible at Norunda.\\
For the data collection procedure, the goal is to cross vary the environmental parameters mentioned above as shown in the preliminary calibration profile in 
%\autoref{fig:calibration_profile}.  
The ranges to be covered for each environmental parameter are:
\begin{itemize}
  \item Humidity: 0-80\%
  \item Temperature: 15-45°C
  \item Pressure: 2.5-3.2 bar
\end{itemize}   
to make sure all possible combination that the K96 might experience during the stratospheric flight are covered.